% !TeX root = ../bnuthesis-example.tex

\chapter{图表和公式的补充示例}

附录是与论文内容密切相关、但编入正文又影响整篇论文编排的条理和逻辑性的资料，例如某些重要的数据表格、计算程序、统计表等，是论文主体的补充内容，可根据需要设置。
本例设置两份附录，展示图、表、公式如何随附录重新编号，正文页码保持连续。


\section{图表示例}

\subsection{图}

附录中的图片示例（图~\ref{fig:appendix-figure}）。

\begin{figure}[htbp]
  \centering
  \includegraphics[width=0.6\linewidth]{example-image-a.pdf}
  \caption{附录中的图片示例}
  \label{fig:appendix-figure}
\end{figure}


\subsection{表格}

附录中的表格示例（表~\ref{tab:appendix-table}）。

\begin{table}[htbp]
  \centering
  \caption{附录中的表格示例}
  \begin{tabular}{ll}
    \toprule
    文件名          & 描述                         \\
    \midrule
    bnuthesis.dtx   & 模板的源文件，包括文档和注释 \\
    bnuthesis.cls   & 模板文件                     \\
    bnuthesis-*.bst & BibTeX 参考文献表样式文件    \\
    bnuthesis-*.bbx & BibLaTeX 参考文献表样式文件  \\
    bnuthesis-*.cbx & BibLaTeX 引用样式文件        \\
    \bottomrule
  \end{tabular}
  \label{tab:appendix-table}
\end{table}


\section{数学公式}

附录中的数学公式示例（公式~\eqref{eq:appendix-equation}）。
\begin{equation}
  \int_a^b f(x)\,\dif x = -\int_b^a f(x)\,\dif x.
  \label{eq:appendix-equation}
\end{equation}
其中，$f$ 是在端点 $a$ 和 $b$ 之间可积的函数。

\chapter{第二份附录的独立编号}

第二份附录用于核对图~\ref{fig:appendix-b-figure}、表~\ref{tab:appendix-b-table} 和式~\eqref{eq:appendix-b-equation} 的序号是否从 B1 开始，避免沿用第一份附录或正文的编号。

\begin{figure}[htbp]
  \centering
  \includegraphics[width=0.6\linewidth]{example-image-b.pdf}
  \caption{第二份附录的图片示例}
  \label{fig:appendix-b-figure}
\end{figure}

\begin{table}[htbp]
  \centering
  \caption{附录编号检查项目}
  \label{tab:appendix-b-table}
  \begin{tabular}{ll}
    \toprule
    项目 & 检查内容 \\
    \midrule
    图与表 & 各附录分别从 1 开始 \\
    公式 & 使用所属附录字母与序号 \\
    页码 & 与正文及其他文后部分连续 \\
    \bottomrule
  \end{tabular}
\end{table}

对正整数 $n$，以 $k$ 为求和指标，有
\begin{equation}
  \sum_{k=1}^{n} k = \frac{n(n+1)}{2}.
  \label{eq:appendix-b-equation}
\end{equation}
